\begin{statementbox}
	\textbf{\textcolor{CStatement}{条件}}

	\begin{description}[style=nextline, leftmargin=2.8em, labelsep=0.5em, itemsep=0.35em]
		\item[\textbf{\textcolor{CStatement}{条件1（椭圆标准形）}}]
		      椭圆方程为
		      \[
			      \frac{x^2}{a^2}+\frac{y^2}{b^2}=1,\qquad a>b>0,
		      \]
		      焦点为 $F(\pm c,0)$，其中
		      \[
			      c^2=a^2-b^2.
		      \]

		\item[\textbf{\textcolor{CStatement}{条件2（过焦直线）}}]
		      直线经过焦点 $F(c,0)$ 或 $F(-c,0)$，并且与椭圆交于两个不同点。
		      若直线不垂直于 $x$ 轴，记其斜率为 $k\in\mathbb R$；若直线垂直于 $x$ 轴，则斜率不存在，需要单独使用通径公式。
	\end{description}

	\textbf{\textcolor{CStatement}{结论}}

	\begin{description}[style=nextline, leftmargin=2.8em, labelsep=0.5em, itemsep=0.45em]
		\item[\textbf{\textcolor{CStatement}{结论一（焦点弦斜率公式）}}]
		      \textbf{\textcolor{CStatement}{直观描述：}}
		      当过焦点的直线不垂直于 $x$ 轴，且斜率为 $k$ 时，焦点弦长为
		      \[
			      |PQ|=\frac{2ab^2(1+k^2)}{a^2k^2+b^2}.
		      \]

		\item[\textbf{\textcolor{CStatement}{结论二（通径公式）}}]
		      \textbf{\textcolor{CStatement}{直观描述：}}
		      当过焦点的直线垂直于长轴时，直线为 $x=c$ 或 $x=-c$，此时斜率不存在。
		      该焦点弦为通径，弦长为
		      \[
			      |PQ|=\frac{2b^2}{a}.
		      \]

		\item[\textbf{\textcolor{CStatement}{补充结论（角度形式）}}]
		      \textbf{\textcolor{CStatement}{直观描述：}}
		      若过焦点的直线与 $x$ 轴的夹角为 $\theta$，则焦点弦长也可写为
		      \[
			      |PQ|=\frac{2ab^2}{a^2\sin^2\theta+b^2\cos^2\theta}.
		      \]
		      这个形式可以直接覆盖垂直情形。
	
        \item[\textbf{\textcolor{CStatement}{补充结论（水平、竖直两条特殊焦点弦）}}]
              \textbf{水平焦点弦：}过任一焦点的水平直线就是长轴 $y=0$，
              端点为 $(-a,0),(a,0)$，因此
              \[
                  \boxed{|PQ|=2a}\qquad\text{（最长）}.
              \]
              \textbf{竖直焦点弦：}过任一焦点的竖直直线 $x=\pm c$
              为通径，其端点纵坐标是 $\pm\frac{b^2}{a}$，因此
              \[
                  \boxed{|PQ|=\frac{2b^2}{a}}\qquad\text{（最短）}.
              \]
              所有过焦点的弦满足
              \[
                  \frac{2b^2}{a}\le |PQ|\le 2a.
              \]

\end{description}
\end{statementbox}